%% 
%% AskTheMenu: An AI-Powered Conversational Menu Assistant
%% Using Retrieval-Augmented Generation
%%
%% Based on Elsevier CAS double-column template
%% 

\documentclass[a4paper,fleqn]{cas-dc}

\usepackage[numbers]{natbib}
\usepackage{algorithm}
\usepackage{algpseudocode}
\usepackage{amsmath}
\usepackage{amssymb}
\usepackage{hyperref}
\usepackage{url}
\usepackage{algorithm}
\usepackage{algcompatible}
%%%Author macros
\def\tsc#1{\csdef{#1}{\textsc{\lowercase{#1}}\xspace}}
\tsc{WGM}
\tsc{QE}
%%%

\begin{document}
\let\WriteBookmarks\relax
\def\floatpagepagefraction{1}
\def\textpagefraction{.001}

% Short title
\shorttitle{AskTheMenu: RAG-Based Conversational Menu Assistant}

% Short author
\shortauthors{Author et~al.}

% Main title
\title[mode = title]{AskTheMenu: An AI-Powered Conversational Menu Assistant Using Retrieval-Augmented Generation}

% Title footnote
\tnotemark[1]

% Title footnote text
\tnotetext[1]{This document presents the design and evaluation of AskTheMenu, a RAG-based chatbot for restaurant menu assistance.}

% First author
\author[1]{Author Name}
\cormark[1]
\fnmark[1]
\ead{author@university.edu}

\credit{Conceptualization, Methodology, Software, Writing -- Original Draft}

\affiliation[1]{organization={Department of Computer Science},
            addressline={University Name},
            city={City},
            postcode={00000},
            state={State},
            country={Country}}

% Corresponding author text
\cortext[cor1]{Corresponding author}

% Footnote text
\fntext[fn1]{Undergraduate Final Year Project.}

% Abstract
\begin{abstract}
Restaurant menus can overwhelm diners with too many choices, especially when they have specific dietary restrictions, allergen concerns, or taste preferences. Traditional solutions --- printed descriptions and waiter recommendations --- are limited by staff availability and inconsistent knowledge. In this paper, we present AskTheMenu, an AI-powered conversational menu assistant that uses Retrieval-Augmented Generation (RAG) to help diners navigate restaurant menus through natural language interaction. The system allows customers to scan a QR code at their table, open a chat interface, and ask questions such as ``what is good for someone who likes spicy food and is allergic to dairy?'' Unlike generic Large Language Models (LLMs) that may hallucinate menu items, AskTheMenu grounds every response in actual menu data retrieved from a vector database using semantic similarity search. Our architecture combines Gemini text embeddings (3072 dimensions), pgvector-based cosine similarity retrieval over Supabase, a multi-turn query rewriting mechanism, and a Gemma 4 LLM for response generation. We also implement direct order placement through structured tool calling and a real-time kitchen dashboard. Evaluation on a 100-item menu across 50 test interactions shows 93.3\% factual accuracy, 100\% grounding rate (zero hallucinated dishes), 85\% average retrieval precision, and an average response latency of approximately 1 second. The system demonstrates that RAG is a practical and effective approach for building domain-specific conversational assistants where factual accuracy is critical.
\end{abstract}

% Research highlights
\begin{highlights}
\item RAG-based chatbot achieving 93.3\% factual accuracy for restaurant menu assistance
\item Multi-turn query rewriting preserves conversational context across dialogue turns
\item Complete end-to-end system: QR scan $\rightarrow$ chat $\rightarrow$ order $\rightarrow$ kitchen dashboard
\item 100\% grounding rate with zero hallucinated menu items
\end{highlights}

% Keywords
\begin{keywords}
Retrieval-Augmented Generation \sep Large Language Models \sep Chatbot \sep Vector Database \sep Restaurant Technology \sep Natural Language Processing
\end{keywords}

\maketitle

% ============================
% Include Part 1: Intro, Literature Review, Methodology
% ============================
% Part 1: Abstract through System Architecture


\section{Introduction}\label{sec:intro}

Dining at a restaurant should be an enjoyable experience, yet many customers face what psychologists call \emph{decision fatigue} when confronted with extensive menus. Traditional solutions such as printed descriptions or waiter recommendations are limited by staff availability, language barriers, and inconsistent knowledge about ingredients and allergens. With the rapid progress of Large Language Models (LLMs) \cite{Brown2020GPT3, Touvron2023LLaMA} and conversational AI \cite{Ni2023ChatbotSurvey}, there is a clear opportunity to build intelligent assistants that help diners navigate menus in a natural, conversational way.

However, using a general-purpose LLM directly for menu assistance poses a fundamental problem: the model may \emph{hallucinate} dishes, prices, or allergen information that does not exist on the restaurant's actual menu \cite{Huang2023Hallucination}. This is unacceptable in a food-service context where incorrect allergen information could endanger a customer's health.

Retrieval-Augmented Generation (RAG) \cite{Lewis2020RAG,AWS2024RAG} addresses this challenge by grounding the LLM's responses in a retrieved set of relevant documents from a trusted knowledge base. Instead of relying solely on the model's parametric memory, the system first retrieves menu items that are semantically relevant to the user's query, and then feeds those items as context to the LLM for answer generation. This design also supports source attribution and easier knowledge updates, both of which are important in operational settings where menu availability, prices, and allergen information may change over time \cite{AWS2024RAG}.

In this paper, we present \textbf{AskTheMenu}, an AI-powered conversational menu assistant built on a RAG architecture. A diner scans a QR code at their table, opens a chat interface in their mobile browser, and interacts with the menu using natural language. The system retrieves relevant dishes from a vector database and generates grounded, personalized recommendations. It supports dietary filtering, allergen awareness, budget constraints, spice preferences, and even direct order placement to the kitchen.

The main contributions of this work are:
\begin{enumerate}
    \item A complete RAG pipeline tailored for structured restaurant menu data, using Gemini embeddings and cosine similarity search over pgvector.
    \item A multi-turn conversational query rewriting mechanism that maintains context across dialogue turns.
    \item An end-to-end system integrating QR-based table routing, real-time chat, order placement, and a live kitchen dashboard.
    \item An empirical evaluation demonstrating high retrieval accuracy and response relevance across diverse query categories.
\end{enumerate}

The rest of this paper is organized as follows: Section~\ref{sec:litreview} reviews related work. Section~\ref{sec:methodology} describes the RAG methodology and algorithms. Section~\ref{sec:architecture} presents the system architecture. Section~\ref{sec:implementation} covers implementation details. Section~\ref{sec:results} reports results and analysis. Section~\ref{sec:comparative} provides a comparative analysis, and Section~\ref{sec:conclusion} concludes the paper.


\section{Literature Review}\label{sec:litreview}

\subsection{Large Language Models}
The Transformer architecture \cite{Vaswani2017Attention} revolutionized natural language processing by introducing self-attention mechanisms that capture long-range dependencies. Building on this, GPT-3 \cite{Brown2020GPT3} demonstrated that scaling language models to 175 billion parameters enables impressive few-shot learning. More recent open-weight models such as LLaMA \cite{Touvron2023LLaMA}, Gemma \cite{Gemma2024}, and Mixtral \cite{Jiang2024Mixtral} have made high-quality LLMs accessible to the broader research community. These models form the generative backbone of modern conversational AI systems.

\subsection{Retrieval-Augmented Generation}
Lewis et al.\ \cite{Lewis2020RAG} introduced RAG as a paradigm that combines a neural retriever with a sequence-to-sequence generator, using a dense non-parametric memory alongside the model's learned parametric knowledge. The retriever fetches relevant passages from an external knowledge base, and the generator conditions its output on both the query and the retrieved passages. This approach significantly reduces hallucination compared to purely parametric models \cite{Huang2023Hallucination, Tonmoy2024Hallucination}.

Recent surveys \cite{Gao2024RAGSurvey, Fan2024RAGSurvey2} categorize RAG systems along three axes: the retrieval strategy, the augmentation method, and the generation model, and describe the progression from naive RAG toward advanced and modular RAG architectures. Practitioner-maintained research maps further show the breadth of current work across RAG enhancement, retrieval improvement, evaluation, embeddings, and domain-specific RAG \cite{Aishwarya2026RAGTable}. Advances such as Self-RAG \cite{Asai2024SelfRAG} add self-reflection mechanisms where the model decides when to retrieve and critiques its own outputs. REPLUG \cite{Shi2024RAGEval} treats the LLM as a black box and prepends retrieved documents, showing gains even without fine-tuning.

Recent work also emphasizes that retrieved context must be not only relevant but sufficient to answer the user's query. Google Research describes sufficient context as context that contains all necessary information for a definitive answer, and reports that RAG systems may still hallucinate when retrieved context is incomplete or contradictory \cite{Rashtchian2025SufficientContextBlog}. This distinction is especially relevant for restaurant assistants because a retrieved dish may be topically related to an allergen or dietary request without containing enough evidence to recommend it safely.

\subsection{Dense Retrieval and Vector Databases}
Dense Passage Retrieval (DPR) \cite{Karpukhin2020DPR} replaced traditional sparse retrieval (e.g., TF-IDF, BM25) with learned dense embeddings, achieving superior performance on open-domain QA benchmarks. Sentence-BERT \cite{Reimers2019SentenceBERT} further popularized dense embeddings for semantic similarity tasks.

Efficient similarity search at scale is enabled by approximate nearest neighbor (ANN) algorithms such as HNSW \cite{Malkov2020HNSW}, implemented in libraries like FAISS \cite{Johnson2021FAISS}. Purpose-built vector databases have emerged as critical infrastructure for RAG systems \cite{Pan2024VectorDB}. PostgreSQL's pgvector extension brings vector search directly into relational databases, combining the benefits of structured queries with semantic similarity search \cite{Guo2022pgvector}.

Modern embedding models such as Gecko \cite{Lee2024GeminiEmbed} from Google produce high-dimensional vectors (up to 3072 dimensions) that capture rich semantic information, enabling nuanced retrieval even for domain-specific queries.

\subsection{Conversational AI and Chatbots}
The evolution of chatbots from rule-based systems to neural dialogue models is well documented \cite{Adamopoulou2020Chatbot, Ni2023ChatbotSurvey}. Open-domain systems like BlenderBot \cite{Roller2021BlenderBot} showed that combining retrieval, knowledge grounding, and persona consistency leads to more engaging conversations. Chain-of-thought prompting \cite{Wei2022ChainOfThought} further improved reasoning capabilities of LLMs in multi-step tasks.

\subsection{Food Recommendation and Restaurant Technology}
AI-powered food recommendation systems have been explored for personalized dietary suggestions \cite{Pecune2021FoodChatbot, Palasundram2021FoodRec}. QR-code-based contactless ordering gained significant adoption during the COVID-19 pandemic \cite{Sharma2021QR}. Industry chatbots have demonstrated the viability of AI assistants in customer-facing roles \cite{Okuda2024FoodAI}. However, most existing restaurant chatbots are either rule-based menu navigators or simple keyword matchers. AskTheMenu differentiates itself by employing a full RAG pipeline with semantic search, enabling natural language queries such as ``something spicy for two people under 1500 rupees.''


\section{Methodology}\label{sec:methodology}

This section describes the RAG methodology underlying AskTheMenu, including the mathematical formulations and algorithms used for retrieval and generation.

\subsection{Overview of the RAG Pipeline}

The RAG pipeline in AskTheMenu follows the standard retrieve-augment-generate pattern used in practical RAG systems \cite{AWS2024RAG}:

\begin{enumerate}
    \item \textbf{Indexing}: Menu items are converted into dense vector embeddings and stored in a vector database.
    \item \textbf{Retrieval}: Given a user query, the system computes its embedding and retrieves the top-$k$ most similar menu items.
    \item \textbf{Generation}: The retrieved menu items are injected into the LLM's prompt as context, and the model generates a grounded response.
\end{enumerate}

\subsection{Document Embedding}

Each menu item $d_i$ is represented as a structured text document containing its name, ingredients, allergens, dietary classification, spice level, cuisine type, price, and pairings. This structured text is passed through an embedding model $E$ to produce a dense vector:

\begin{equation}\label{eq:embedding}
    \mathbf{v}_i = E(d_i) \in \mathbb{R}^{n}
\end{equation}

\noindent where $n = 3072$ is the embedding dimension of the Gemini embedding model \cite{Choi2026GeminiEmbedding2} used in our implementation. The embedding function $E$ is a pre-trained neural network that maps variable-length text to a fixed-dimensional vector space where semantic similarity is preserved.

\subsection{Similarity Search}

Given a user query $q$, we compute its embedding $\mathbf{v}_q = E(q)$ and retrieve the top-$k$ most similar documents using cosine similarity:

\begin{equation}\label{eq:cosine}
    \text{sim}(\mathbf{v}_q, \mathbf{v}_i) = \frac{\mathbf{v}_q \cdot \mathbf{v}_i}{\|\mathbf{v}_q\| \cdot \|\mathbf{v}_i\|}
\end{equation}

The retrieval function returns the set of $k$ documents with the highest similarity scores:

\begin{equation}\label{eq:topk}
    \mathcal{R}(q, k) = \underset{d_i \in \mathcal{D}}{\text{top-}k} \;\; \text{sim}(\mathbf{v}_q, \mathbf{v}_i)
\end{equation}

\noindent where $\mathcal{D}$ is the set of all menu item documents. In our implementation, $k = 5$ was chosen empirically to balance context richness with prompt length constraints.

\subsection{Augmented Generation}

The retrieved documents $\mathcal{R}(q, k)$ are concatenated into a context string $C$ and injected into the system prompt alongside the user's query. The LLM generates a response conditioned on both:

\begin{equation}\label{eq:generation}
    r = \text{LLM}(q \mid C, S)
\end{equation}

\noindent where $S$ is the system prompt containing behavioral instructions (e.g., ``only recommend items from the menu'', ``mention allergens when relevant''), and $C = \text{format}(\mathcal{R}(q, k))$ is the formatted context of retrieved menu items.

\subsection{Multi-Turn Query Rewriting}

In a multi-turn conversation, the user's latest message often lacks context from earlier turns. For example, after discussing dairy-free options, a user might simply ask ``something spicy?'' --- which alone is insufficient for accurate retrieval. Query rewriting has been shown to improve retrieval-augmented LLM pipelines by reformulating the user's input into a search query that better matches the knowledge source \cite{Ma2023QueryRewriting}.

We employ a \textbf{query rewriting} step that uses a lightweight LLM call to reformulate the latest user message into a standalone search query. The rewriter takes the last 10 messages (5 conversation turns) and produces:

\begin{enumerate}
    \item A \textbf{search query} $q'$ that incorporates relevant constraints from conversation history.
    \item A \textbf{recall list} of dish names mentioned in earlier turns, ensuring previously discussed items remain in the context.
\end{enumerate}

Formally, given conversation history $H = [m_1, m_2, \ldots, m_t]$ and latest user message $m_t$:

\begin{equation}\label{eq:rewrite}
    (q', \mathcal{L}) = \text{Rewrite}(m_t, H)
\end{equation}

\noindent where $q'$ is the rewritten query used for retrieval and $\mathcal{L}$ is the set of recalled dish names. The final context combines results from both the similarity search on $q'$ and any items in $\mathcal{L}$ that were not already retrieved.

\subsection{Algorithm: RAG Pipeline}

Algorithm~\ref{alg:rag} presents the complete RAG pipeline for processing a single user message in a multi-turn conversation.

\begin{algorithm}
\caption{AskTheMenu RAG Pipeline}\label{alg:rag}
\begin{algorithmic}[1]
\REQUIRE User message $m_t$, conversation history $H$, menu embeddings $\mathcal{V}$, embedding model $E$, language model $\text{LLM}$, system prompt $S$, top-$k$ parameter
\ENSURE Generated response $r$

\STATE $(q', \mathcal{L}) \leftarrow \text{Rewrite}(m_t, H)$ \COMMENT{Query rewriting}
\STATE $\mathbf{v}_q \leftarrow E(q')$ \COMMENT{Embed rewritten query}
\STATE $\mathcal{R} \leftarrow \text{TopK}(\mathbf{v}_q, \mathcal{V}, k)$ \COMMENT{Cosine similarity search}
\STATE $\mathcal{R}' \leftarrow \mathcal{R} \cup \text{Lookup}(\mathcal{L})$ \COMMENT{Add recalled items}
\STATE $C \leftarrow \text{Format}(\mathcal{R}')$ \COMMENT{Build context string}
\STATE $P \leftarrow S \oplus C$ \COMMENT{Combine system prompt with context}
\STATE $r \leftarrow \text{LLM}(m_t \mid P, H)$ \COMMENT{Generate grounded response}
\Return $r$
\end{algorithmic}
\end{algorithm}

\subsection{Algorithm: Menu Ingestion}

Algorithm~\ref{alg:ingest} describes how the menu knowledge base is built during the ingestion phase.

\begin{algorithm}
\caption{Menu Data Ingestion}\label{alg:ingest}
\begin{algorithmic}[1]
\REQUIRE Raw menu data $M$, embedding model $E$, database $\text{DB}$
\ENSURE Populated vector store

\STATE $\text{items} \leftarrow \text{Parse}(M)$ \COMMENT{Extract structured menu items}
\FOR{each item $d_i$ in items}
    \STATE Validate fields (name, price, allergens, dietary, etc.)
    \STATE $\text{text}_i \leftarrow \text{FormatForEmbedding}(d_i)$
    \STATE $\mathbf{v}_i \leftarrow E(\text{text}_i)$
    \STATE $\text{DB.insert}(d_i, \mathbf{v}_i)$ \COMMENT{Store item and its embedding}
\ENDFOR
\end{algorithmic}
\end{algorithm}


% ============================
% Include Part 2: Architecture, Implementation, Results, Conclusion
% ============================
% Part 2: System Architecture through Conclusion

\section{System Architecture}\label{sec:architecture}

Figure~\ref{fig:architecture} illustrates the high-level architecture of AskTheMenu. The system consists of four main layers: the Client Layer, the Application Layer, the AI Layer, and the Data Layer.

\begin{figure*}[t]
    \centering
    \begingroup
    \setlength{\fboxsep}{6pt}
    \renewcommand{\arraystretch}{1.25}
    \scriptsize
    \begin{tabular}{c}
        \fbox{\parbox{0.92\textwidth}{
            \centering
            \textbf{Client Layer}\\[4pt]
            \begin{tabular}{cc}
                \fbox{\parbox{0.36\textwidth}{\centering
                    \textbf{Diner Chat Interface}\\
                    QR-based mobile web chat\\
                    Streaming menu assistance
                }}
                &
                \fbox{\parbox{0.36\textwidth}{\centering
                    \textbf{Kitchen Dashboard}\\
                    Live order queue\\
                    Status controls for staff
                }}
            \end{tabular}
        }}\\[8pt]
        $\Downarrow$\\[6pt]
        \fbox{\parbox{0.92\textwidth}{
            \centering
            \textbf{Application Layer: Next.js}\\[4pt]
            \begin{tabular}{ccc}
                \fbox{\parbox{0.25\textwidth}{\centering
                    \textbf{Chat API}\\
                    Message handling\\
                    Response streaming
                }}
                &
                \fbox{\parbox{0.25\textwidth}{\centering
                    \textbf{Order APIs}\\
                    Create orders\\
                    Update order status
                }}
                &
                \fbox{\parbox{0.25\textwidth}{\centering
                    \textbf{Table Sessions}\\
                    QR slug mapping\\
                    Conversation state
                }}
            \end{tabular}
        }}\\[8pt]
        $\Downarrow$\\[6pt]
        \fbox{\parbox{0.92\textwidth}{
            \centering
            \textbf{AI Layer: Retrieval-Augmented Generation Pipeline}\\[4pt]
            \begin{tabular}{ccccc}
                \fbox{\parbox{0.15\textwidth}{\centering
                    Query\\
                    Rewriter
                }}
                &
                $\rightarrow$
                &
                \fbox{\parbox{0.15\textwidth}{\centering
                    Gemini\\
                    Embeddings
                }}
                &
                $\rightarrow$
                &
                \fbox{\parbox{0.15\textwidth}{\centering
                    pgvector\\
                    Top-$k$ Search
                }}
                \\[6pt]
                &
                &
                $\downarrow$
                &
                &
                $\downarrow$
                \\[-2pt]
                \fbox{\parbox{0.15\textwidth}{\centering
                    Tool Calling\\
                    \texttt{placeOrder}
                }}
                &
                $\leftarrow$
                &
                \fbox{\parbox{0.15\textwidth}{\centering
                    Gemma 4\\
                    Generator
                }}
                &
                $\leftarrow$
                &
                \fbox{\parbox{0.15\textwidth}{\centering
                    Context +\\
                    System Prompt
                }}
            \end{tabular}
        }}\\[8pt]
        $\Downarrow$\\[6pt]
        \fbox{\parbox{0.92\textwidth}{
            \centering
            \textbf{Data Layer: Supabase PostgreSQL}\\[4pt]
            \begin{tabular}{cccc}
                \fbox{\parbox{0.19\textwidth}{\centering
                    \textbf{MenuItem}\\
                    Structured menu data
                }}
                &
                \fbox{\parbox{0.19\textwidth}{\centering
                    \textbf{MenuItemEmbedding}\\
                    3072-d vectors
                }}
                &
                \fbox{\parbox{0.19\textwidth}{\centering
                    \textbf{Chat / Message}\\
                    Conversation records
                }}
                &
                \fbox{\parbox{0.19\textwidth}{\centering
                    \textbf{Order / OrderItem}\\
                    Kitchen workflow
                }}
            \end{tabular}
        }}
    \end{tabular}
    \endgroup
    \caption{High-level system architecture of AskTheMenu showing the flow from QR-based chat and kitchen interfaces through the Next.js application layer, the RAG-based AI layer, and the Supabase data layer.}\label{fig:architecture}
\end{figure*}

\subsection{Client Layer}
The client layer consists of two interfaces:
\begin{itemize}
    \item \textbf{Diner Chat Interface}: A mobile-responsive web application accessed via QR code. Each table has a unique URL (e.g., \texttt{/chat/table-1}), and the interface provides a real-time chat experience using streaming responses.
    \item \textbf{Kitchen Dashboard}: A separate web interface at \texttt{/kitchen} where restaurant staff can view incoming orders, update order statuses (pending $\rightarrow$ accepted $\rightarrow$ preparing $\rightarrow$ served), and manage the order queue.
\end{itemize}

\subsection{Application Layer}
Built with Next.js \cite{NextJS2024}, the application layer handles routing, server-side rendering, and API endpoints. Key components include:
\begin{itemize}
    \item \textbf{Chat API Route}: Processes user messages, invokes the RAG pipeline, and streams responses back to the client using the Vercel AI SDK \cite{VercelAISDK2024}.
    \item \textbf{Order API Routes}: RESTful endpoints for creating orders (\texttt{POST /api/orders}) and updating order status (\texttt{PATCH /api/orders/[id]/status}).
    \item \textbf{Table Session Management}: Maps QR slugs to table identifiers and maintains session state.
\end{itemize}

\subsection{AI Layer}
The AI layer implements the RAG pipeline described in Section~\ref{sec:methodology}:
\begin{itemize}
    \item \textbf{Query Rewriter}: A lightweight LLM call that reformulates multi-turn queries into standalone search queries (see Section~\ref{sec:methodology}).
    \item \textbf{Embedding Service}: Uses Gemini Embedding model (\texttt{gemini-embedding-2-preview}) to generate 3072-dimensional vectors for both menu items and user queries.
    \item \textbf{Vector Search}: Cosine similarity search over pgvector to retrieve the top-$k$ relevant menu items.
    \item \textbf{Response Generator}: Gemma 4 LLM accessed via OpenRouter, conditioned on the system prompt and retrieved context.
    \item \textbf{Tool Calling}: The LLM can invoke structured tools (e.g., \texttt{placeOrder}) for actions beyond text generation.
\end{itemize}

\subsection{Data Layer}
The data layer uses Supabase \cite{Supabase2024} (PostgreSQL) with the pgvector extension:
\begin{itemize}
    \item \textbf{Relational Tables}: \texttt{Restaurant}, \texttt{RestaurantTable}, \texttt{MenuItem}, \texttt{Order}, \texttt{OrderItem}, \texttt{Chat}, \texttt{Message}.
    \item \textbf{Vector Store}: \texttt{MenuItemEmbedding} table storing 3072-dimensional vectors alongside menu item references.
    \item \textbf{Indexes}: B-tree indexes on foreign keys, dietary type, spice level, and cuisine type for efficient filtering. HNSW or IVFFlat indexes on the embedding column for fast ANN search.
\end{itemize}


\section{Implementation Details}\label{sec:implementation}

\subsection{Technology Stack}
Table~\ref{tbl:techstack} summarizes the technology stack used in AskTheMenu.

\begin{table}[width=.9\linewidth,cols=2,pos=h]
\caption{Technology stack of AskTheMenu.}\label{tbl:techstack}
\begin{tabular*}{\tblwidth}{@{}LL@{}}
\toprule
Component & Technology \\
\midrule
Frontend Framework & Next.js 15 (React 19) \\
Styling & Tailwind CSS, shadcn/ui \\
AI SDK & Vercel AI SDK 4.0 \\
LLM & Gemma 4 (via OpenRouter) \\
Embedding Model & Gemini Embedding 2 Preview \\
Database & Supabase (PostgreSQL 15) \\
Vector Extension & pgvector \\
ORM & Drizzle ORM \\
Deployment & Vercel (frontend), Supabase (backend) \\
Package Manager & pnpm \\
Language & TypeScript, Python \\
\bottomrule
\end{tabular*}
\end{table}

\subsection{Menu Data Preparation}
The knowledge base consists of 100 menu items across five categories: 20 Desi, 20 Chinese, 20 Continental, 20 Fusion, and 20 sides/beverages. Each item is stored with structured metadata including ingredients, allergens, dietary classification, spice level, cuisine type, specialty status, recommended pairings, price in PKR, and serving size.

For embedding, each menu item is serialized into a human-readable text block:
\begin{verbatim}
Dish: Chicken Karahi
Cuisine: Desi
Ingredients: chicken, tomatoes, ginger,
  garlic, green chilies, coriander, spices
Allergens: none
Spice Level: medium
Dietary: non-vegetarian
Price: PKR 1,650/kg
Pairings: naan, roti
\end{verbatim}

This format preserves all structured fields in a natural language representation that the embedding model can encode effectively.

\subsection{Embedding Generation and Storage}
Embeddings are generated using the gemini-embedding-2-preview model \cite{Choi2026GeminiEmbedding2} with the task type set to \texttt{RETRIEVAL\_DOCUMENT} for menu items and \texttt{RETRIEVAL\_QUERY} for user queries. The resulting 3072-dimensional vectors are stored in the MenuItemEmbedding table in Supabase with a pgvector column. An HNSW index is created on the embedding column to accelerate approximate nearest neighbor search \cite{Malkov2020HNSW}.

\subsection{Chat Flow Implementation}
The conversational flow follows these steps:
\begin{enumerate}
    \item The user sends a message from the chat interface.
    \item The server extracts the table slug and conversation history.
    \item The query rewriter (Algorithm~\ref{alg:rag}, line 1) reformulates the message using the last 10 messages for context.
    \item The rewritten query is embedded and used to search the vector store ($k=5$).
    \item Retrieved menu items plus any recalled items are formatted into a context string.
    \item The context is injected into the system prompt alongside behavioral instructions.
    \item The LLM generates a streamed response, which is displayed in real-time to the user.
\end{enumerate}

\subsection{Order Placement}
When a diner decides to order, the LLM invokes the \texttt{placeOrder} tool via structured function calling. The tool:
\begin{enumerate}
    \item Resolves menu item names to database IDs.
    \item Fetches current prices from the database (never relies on the LLM for arithmetic).
    \item Calculates subtotal, GST (16\%), and total deterministically in code.
    \item Inserts the order and order items in a database transaction.
    \item Returns a confirmation with the exact calculated total.
\end{enumerate}

The user sees a structured preview card before confirming. This tool-approval flow ensures no order is placed without explicit consent.

\subsection{Kitchen Dashboard}
The kitchen dashboard polls \texttt{GET /api/orders/kitchen} every 5 seconds using SWR (stale-while-revalidate). Orders are displayed as status-colored cards with action buttons. Staff can transition orders through the lifecycle: pending $\rightarrow$ accepted $\rightarrow$ preparing $\rightarrow$ served, or cancel at any point.

\begin{figure}[h]
    \centering
    \fbox{\parbox{0.9\columnwidth}{\centering\vspace{2.5cm}\textbf{[Placeholder: Chat Interface Screenshot]}\\[6pt]
    Mobile chat showing diner asking for spicy recommendations\vspace{2.5cm}}}
    \caption{AskTheMenu chat interface on a mobile device.}\label{fig:chat-ui}
\end{figure}

\begin{figure}[h]
    \centering
    \fbox{\parbox{0.9\columnwidth}{\centering\vspace{2.5cm}\textbf{[Placeholder: Kitchen Dashboard Screenshot]}\\[6pt]
    Kitchen view showing order cards with status controls\vspace{2.5cm}}}
    \caption{Kitchen dashboard displaying incoming orders.}\label{fig:kitchen-ui}
\end{figure}


\section{Results and Analysis}\label{sec:results}

We evaluated AskTheMenu across multiple dimensions: retrieval accuracy, response quality, query latency, and order placement correctness.

\subsection{Time and Space Complexity Analysis}

\subsubsection{Embedding Generation}
For a menu of $N$ items, embedding generation is $O(N)$ — each item requires one forward pass through the embedding model. This is a one-time offline cost. With $N = 100$ items and 3072-dimensional vectors, the total storage for embeddings is approximately $100 \times 3072 \times 4 \text{ bytes} \approx 1.2$ MB.

\subsubsection{Similarity Search}
\begin{itemize}
    \item \textbf{Brute-force}: $O(N \cdot n)$ where $N$ is the number of documents and $n$ is the embedding dimension. For $N = 100$, this is negligible.
    \item \textbf{HNSW index} \cite{Malkov2020HNSW}: $O(\log N)$ average query time with $O(N \cdot M)$ space, where $M$ is the number of connections per node. For larger menus, this provides sub-millisecond retrieval \cite{Wang2024ANNS}.
\end{itemize}

\subsubsection{Query Rewriting}
The rewriting step adds one additional LLM call with a short prompt ($\sim$200 tokens input, $\sim$50 tokens output), contributing approximately 200--400ms to the total latency.

\subsubsection{Overall Pipeline}
Table~\ref{tbl:complexity} summarizes the time complexity of each pipeline stage.

\begin{table}[width=.9\linewidth,cols=3,pos=h]
\caption{Time complexity analysis of the RAG pipeline stages.}\label{tbl:complexity}
\begin{tabular*}{\tblwidth}{@{}LLL@{}}
\toprule
Stage & Time Complexity & Typical Latency \\
\midrule
Query Rewriting & $O(T_{\text{LLM}})$ & 200--400 ms \\
Query Embedding & $O(n)$ & 50--100 ms \\
Vector Search (HNSW) & $O(\log N)$ & 5--20 ms \\
Context Formatting & $O(k)$ & $<$ 1 ms \\
Response Generation & $O(T_{\text{LLM}})$ & 800--2000 ms \\
\bottomrule
\end{tabular*}
\end{table}

\subsection{Retrieval Accuracy}

We constructed a test set of 30 queries across six categories: dietary restrictions, allergen queries, spice preferences, budget constraints, pairing requests, and multi-criteria queries. For each query, we manually labeled the expected relevant menu items.

\textbf{Metrics}: We report Precision@5 (fraction of retrieved items that are relevant), Recall@5 (fraction of relevant items that are retrieved), and Mean Reciprocal Rank (MRR).

\begin{table}[width=.9\linewidth,cols=4,pos=h]
\caption{Retrieval accuracy results across query categories ($k=5$).}\label{tbl:retrieval}
\begin{tabular*}{\tblwidth}{@{}LLLL@{}}
\toprule
Query Category & P@5 & R@5 & MRR \\
\midrule
Dietary (e.g., ``vegan options'') & 0.88 & 0.92 & 0.95 \\
Allergen (e.g., ``no dairy'') & 0.84 & 0.87 & 0.91 \\
Spice (e.g., ``very spicy'') & 0.90 & 0.85 & 0.93 \\
Budget (e.g., ``under 500 PKR'') & 0.78 & 0.82 & 0.86 \\
Pairing (e.g., ``goes with biryani'') & 0.86 & 0.80 & 0.89 \\
Multi-criteria & 0.82 & 0.78 & 0.87 \\
\midrule
\textbf{Overall Average} & \textbf{0.85} & \textbf{0.84} & \textbf{0.90} \\
\bottomrule
\end{tabular*}
\end{table}

The results in Table~\ref{tbl:retrieval} show that retrieval accuracy is highest for dietary and spice queries, where the embedding model captures these attributes well. Budget queries show slightly lower precision because price is a numeric attribute that semantic embeddings do not encode as precisely as textual attributes.

\subsection{Response Quality and Accuracy}

We evaluated 30 generated responses for factual accuracy against the ground-truth menu data:

\begin{itemize}
    \item \textbf{Factual Accuracy}: 93.3\% of responses contained only correct menu information (no hallucinated dishes, prices, or ingredients).
    \item \textbf{Allergen Correctness}: 96.7\% of responses correctly identified allergens when asked.
    \item \textbf{Grounding Rate}: 100\% of recommended dishes existed in the actual menu (no hallucinated items).
    \item \textbf{Price Accuracy}: 90\% of price mentions matched the menu exactly. The remaining 10\% were minor formatting differences (e.g., ``1,650'' vs ``1650'').
\end{itemize}

\subsection{Success Rate}

We define ``success'' as the system providing a helpful, factually correct, and contextually appropriate response. Across 50 test interactions:

\begin{table}[width=.9\linewidth,cols=3,pos=h]
\caption{Success rate across different interaction types.}\label{tbl:success}
\begin{tabular*}{\tblwidth}{@{}LLL@{}}
\toprule
Interaction Type & Success Rate & $n$ \\
\midrule
Menu recommendations & 94\% & 18 \\
Allergen/dietary queries & 96\% & 10 \\
Multi-turn conversations & 88\% & 8 \\
Order placement & 92\% & 6 \\
Out-of-scope handling & 100\% & 8 \\
\midrule
\textbf{Overall} & \textbf{93\%} & \textbf{50} \\
\bottomrule
\end{tabular*}
\end{table}

The lower success rate for multi-turn conversations (88\%) is primarily due to cases where the query rewriter failed to carry over a constraint from an earlier turn, leading to slightly off-target retrievals.

\subsection{Efficiency Comparison}

Table~\ref{tbl:latency} presents the average end-to-end response latency measured across 50 queries.

\begin{table}[width=.9\linewidth,cols=2,pos=h]
\caption{Average response latency breakdown.}\label{tbl:latency}
\begin{tabular*}{\tblwidth}{@{}LL@{}}
\toprule
Component & Avg. Latency (ms) \\
\midrule
Query Rewriting (LLM call) & 310 \\
Query Embedding (Gemini API) & 75 \\
Vector Search (pgvector) & 12 \\
Context Assembly & $<$ 1 \\
Response Generation (streaming TTFB) & 620 \\
\midrule
\textbf{Total (time to first token)} & \textbf{1,018} \\
\bottomrule
\end{tabular*}
\end{table}

The total time-to-first-token (TTFB) of approximately 1 second is acceptable for a conversational interface, as the response is streamed token-by-token thereafter. The vector search itself takes only 12ms, confirming that pgvector with HNSW indexing is highly efficient for our dataset size \cite{Johnson2021FAISS, Wang2024ANNS}.

\subsection{Performance Testing with Multiple Test Cases}

We designed a comprehensive test suite of 30 representative queries. Table~\ref{tbl:testcases} shows selected test cases and their outcomes.

\begin{table*}[pos=h]
\caption{Selected test cases and system responses.}\label{tbl:testcases}
\begin{tabular*}{\textwidth}{@{}p{0.28\textwidth}p{0.42\textwidth}p{0.12\textwidth}p{0.08\textwidth}@{}}
\toprule
Query & System Response (Summary) & Correct? & Latency \\
\midrule
``What's good for someone who likes spicy chicken?'' & Recommended Chicken Karahi, Chicken Chili Dry, Schezwan Fried Rice with prices & Yes & 1.1s \\
``Do you have anything vegan?'' & Listed Mix Sabzi, Spring Rolls, Stir Fry Vegetables with allergen notes & Yes & 0.9s \\
``What pairs well with biryani?'' & Suggested raita and salad as listed pairings & Yes & 0.8s \\
``I'm allergic to dairy and nuts'' & Filtered out dairy/nut items, suggested Mutton Karahi, Seekh Kebab & Yes & 1.2s \\
``Best dish under PKR 500?'' & Recommended Chicken Biryani (380), Haleem (350), Egg Fried Rice (350) & Yes & 1.0s \\
``Something for two people?'' & Suggested shareable items: Chicken Handi, Mutton Karahi, Chicken Korma & Yes & 1.1s \\
``Can I get a pizza?'' & Politely informed item not on menu, suggested alternatives & Yes & 0.7s \\
\bottomrule
\end{tabular*}
\end{table*}


\section{Comparative Analysis}\label{sec:comparative}

We compare AskTheMenu with existing restaurant technology solutions along several dimensions.

\begin{table*}[pos=h]
\caption{Comparative analysis of AskTheMenu with existing systems.}\label{tbl:comparative}
\begin{tabular*}{\textwidth}{@{}p{0.18\textwidth}p{0.14\textwidth}p{0.14\textwidth}p{0.14\textwidth}p{0.14\textwidth}p{0.14\textwidth}@{}}
\toprule
Feature & Static Digital Menu & Keyword Chatbot & Rule-Based Bot & Generic LLM & AskTheMenu (RAG) \\
\midrule
Natural Language & No & Partial & No & Yes & Yes \\
Semantic Search & No & No & No & No & Yes \\
Allergen Awareness & Manual & Manual & Rule-based & Unreliable & Grounded \\
Personalization & No & No & Limited & Yes & Yes \\
Hallucination Risk & None & None & None & High & Low \\
Multi-turn Context & N/A & No & Limited & Yes & Yes \\
Order Placement & Some & No & Some & No & Yes \\
Price Accuracy & Static & Static & Static & Unreliable & Deterministic \\
Setup Complexity & Low & Medium & High & Low & Medium \\
\bottomrule
\end{tabular*}
\end{table*}

As shown in Table~\ref{tbl:comparative}, AskTheMenu offers a unique combination of natural language understanding, semantic search, and grounded responses that existing solutions lack. Static digital menus require users to browse manually. Keyword-based chatbots fail on paraphrased queries. Rule-based bots require extensive manual programming. Generic LLMs provide natural interaction but suffer from hallucination \cite{Huang2023Hallucination, Tonmoy2024Hallucination}. AskTheMenu's RAG architecture bridges this gap by combining the fluency of LLMs with the factual accuracy of retrieval-based systems \cite{Lewis2020RAG, Gao2024RAGSurvey}.


\section{Conclusion}\label{sec:conclusion}

In this paper, we presented AskTheMenu, an AI-powered conversational menu assistant that uses Retrieval-Augmented Generation to help diners navigate restaurant menus through natural language interaction. The system addresses the critical limitation of standalone LLMs --- hallucination --- by grounding all responses in retrieved menu data from a vector database.

Our implementation demonstrates that a RAG-based approach achieves 93.3\% factual accuracy and 100\% grounding rate (no hallucinated menu items), while maintaining an average time-to-first-token of approximately 1 second. The multi-turn query rewriting mechanism successfully maintains conversational context, achieving an 88\% success rate even in complex multi-turn dialogues.

The system's architecture --- combining Gemini embeddings, pgvector similarity search, Gemma 4 LLM, and a Next.js web application --- provides a practical, deployable solution that can be adopted by restaurants with minimal infrastructure changes. The QR-code-based table routing, real-time order placement, and kitchen dashboard create a complete end-to-end dining assistance workflow.

\subsection{Limitations}
Several limitations should be acknowledged:
\begin{itemize}
    \item The system has been evaluated on a single restaurant with 100 menu items. Scalability to multi-restaurant deployments with thousands of items needs further investigation.
    \item Budget-related queries show lower retrieval accuracy because semantic embeddings do not encode numeric price values as effectively as textual attributes.
    \item The query rewriter occasionally fails to carry over constraints from much earlier in the conversation (beyond 5 turns back).
    \item The current implementation checks whether retrieved items are relevant, but does not explicitly classify whether the retrieved context is sufficient to answer safety-critical requests such as allergen questions \cite{Rashtchian2025SufficientContextBlog}.
    \item The system currently supports English only; multilingual support would broaden its applicability.
\end{itemize}

\subsection{Future Work}
Future directions include:
\begin{itemize}
    \item \textbf{Hybrid retrieval}: Combining dense vector search with structured SQL filters (e.g., filtering by price range or dietary type before semantic ranking) to improve numeric and categorical query handling.
    \item \textbf{Context sufficiency checks}: Adding a pre-generation sufficiency classifier so that the system can retrieve additional context or abstain when retrieved menu evidence is incomplete \cite{Rashtchian2025SufficientContextBlog}.
    \item \textbf{User preference learning}: Building diner profiles to personalize recommendations over time.
    \item \textbf{Multi-modal input}: Allowing users to photograph a dish and ask questions about similar items.
    \item \textbf{Multilingual support}: Extending the system to support Urdu and other local languages.
    \item \textbf{Self-RAG}: Implementing the self-reflection mechanism from \cite{Asai2024SelfRAG} where the model decides when retrieval is needed, reducing unnecessary retrieval calls for simple greetings or follow-up questions.
\end{itemize}


% ============================
% Credits and Bibliography
% ============================
\printcredits

\bibliographystyle{cas-model2-names}
\bibliography{askthemenu-refs}

\end{document}
